%============================================================
% BAB 3: LANDASAN TEORI
%   Flat list per teknologi/konsep, masing-masing jadi section.
%   Sub-section untuk sub-topik teknologi jika diperlukan.
%   Sesuaikan dengan teknologi yang digunakan pada proyek Anda.
%============================================================
\chapter{LANDASAN TEORI}

Bab ini memuat teori-teori dan teknologi yang menjadi landasan dalam pelaksanaan
kerja praktik di \NamaPerusahaan. Teori yang dibahas meliputi konsep dasar,
arsitektur, kelebihan, dan implementasi umum dari teknologi-teknologi yang
digunakan dalam proyek selama kerja praktik berlangsung.

% -- Sesuaikan section berikut dengan teknologi yang digunakan ----------------

\section{[TEKNOLOGI/KONSEP PERTAMA]}

Jelaskan konsep dasar teknologi atau framework pertama yang digunakan.
Sertakan definisi, tujuan penggunaan, dan bagaimana teknologi ini
berkontribusi pada proyek selama kerja praktik.

% Tambahkan sub-section jika teknologi ini memiliki sub-topik penting.
\subsection{[Sub-topik jika diperlukan]}

Jelaskan sub-topik pertama dari teknologi ini.

\section{[TEKNOLOGI/KONSEP KEDUA]}

Jelaskan konsep dasar teknologi atau framework kedua yang digunakan.
Sertakan definisi, tujuan penggunaan, jika diperlukan.

\section{[TEKNOLOGI/KONSEP KETIGA]}

Jelaskan konsep dasar teknologi atau framework ketiga yang digunakan.

\subsection{[Sub-topik 3.x.1]}

Jelaskan sub-topik pertama.

\subsection{[Sub-topik 3.x.2]}

Jelaskan sub-topik kedua.

% Tambahkan section lebih lanjut sesuai jumlah teknologi yang digunakan.
